\subsection{Выбор программных средств для мобильного приложения}

Выбор технологического стека определялся несколькими требованиями: необходимостью обеспечить работу приложения одновременно на платформах Android и iOS из единой кодовой базы, интеграцией с существующим серверным программным интерфейсом на ASP.NET Core, доступом к аппаратным возможностям устройства (GPS-модуль, камера), а также преемственностью с технологическим стеком диспетчерской веб-панели, реализованной на React и TypeScript.

В качестве основной платформы выбран React Native с языком TypeScript. Главным доводом стала кроссплатформенность и преемственность с веб-панелью диспетчера: подходы к структуре компонентов, управлению состоянием и типизации остаются общими, что снижает порог вхождения и обеспечивает единообразие структурных решений.

Для ускорения разработки использован инструментарий Expo, предоставляющий готовые модули для доступа к камере и геолокации, а также среду Expo Go для тестирования без нативной сборки. Остальные компоненты стека: навигация~--- expo-router с файловой маршрутизацией, доступ к камере~--- expo-image-picker, геолокация~--- expo-location, защищённое хранение токенов~--- expo-secure-store, карта~--- react-native-maps (Google Maps на Android, Apple Maps на iOS), HTTP-взаимодействие~--- Axios с поддержкой JWT-аутентификации.

\subsection{Реализация пользовательского интерфейса мобильного приложения}

Пользовательский интерфейс мобильного приложения реализован в соответствии с макетами, разработанными на этапе проектирования. Ниже описаны реализованные экраны.

\textbf{Экран авторизации} (рисунок~\ref{fig:impl-auth}) является первым экраном для неаутентифицированного пользователя. Содержит поля ввода электронной почты и пароля с возможностью переключения видимости символов, ссылку для восстановления пароля, кнопку <<Войти>>, альтернативный сценарий входа через внешний аккаунт и переход к регистрации. После успешной аутентификации пользователь перенаправляется на основной экран.

\begin{figure}[H]
    \centering
    \setlength{\fboxrule}{0.4pt}
    \setlength{\fboxsep}{0pt}
    \fbox{\includegraphics[width=0.3\linewidth]{mobile-auth-screen.png}}
    \caption{Экран авторизации мобильного приложения}
    \label{fig:impl-auth}
\end{figure}

\textbf{Экран карты} (рисунок~\ref{fig:impl-map}) является центральным экраном приложения. Реализация основана на компоненте MapView из библиотеки react-native-maps. Карта отображает маркеры обращений, цвет которых определяется текущим статусом. Реализованы панель фильтров по категориям, кнопка определения текущего местоположения и плавающая кнопка создания обращения. При нажатии на маркер отображается нижняя информационная карточка.

\begin{figure}[H]
    \centering
    \setlength{\fboxrule}{0.4pt}
    \setlength{\fboxsep}{0pt}
    \fbox{\includegraphics[width=0.35\linewidth]{mobile-map-screen.png}}
    \caption{Экран карты мобильного приложения}
    \label{fig:impl-map}
\end{figure}

\textbf{Экран создания обращения} (рисунок~\ref{fig:impl-create}) предназначен для фиксации проблемы. Реализован в виде модального окна с формой в ScrollView. Основные компоненты: выпадающий список выбора категории, текстовое поле описания, область прикрепления фотографий с использованием expo-image-picker, интерактивная карта для указания местоположения. При обнаружении аналогичного обращения в радиусе 50~м (механизм дедупликации) отображается информационное предупреждение с предложением подтвердить существующее обращение.

\begin{figure}[H]
    \centering
    \setlength{\fboxrule}{0.4pt}
    \setlength{\fboxsep}{0pt}
    \fbox{\includegraphics[width=0.25\linewidth]{mobile-create-report-screen.png}}
    \caption{Экран создания обращения}
    \label{fig:impl-create}
\end{figure}

\textbf{Экран списка обращений} (рисунок~\ref{fig:impl-list}) предназначен для просмотра всех обращений текущего пользователя. Реализован с использованием компонента FlatList с оптимизациями для производительности. Экран содержит панель вкладок для фильтрации по статусу (<<Все>>, <<Активные>>, <<Завершённые>>), поисковую строку с дебаунсом и список карточек. Каждая карточка отображает миниатюру фотографии, категорию и текущий статус. Реализованы индикатор загрузки и механизм обновления списка при подтягивании (RefreshControl).

\begin{figure}[H]
    \centering
    \setlength{\fboxrule}{0.4pt}
    \setlength{\fboxsep}{0pt}
    \fbox{\includegraphics[width=0.25\linewidth]{mobile-reports-list-screen.png}}
    \caption{Экран списка обращений пользователя}
    \label{fig:impl-list}
\end{figure}

\textbf{Экран карточки обращения} (рисунок~\ref{fig:impl-detail}) отображает подробную информацию: фотографию в полную ширину, категорию, адрес, дату создания, описание, счётчик подтверждений от других пользователей и хронологию изменения статусов. При наличии прав доступа отображается кнопка редактирования.

\begin{figure}[H]
    \centering
    \setlength{\fboxrule}{0.4pt}
    \setlength{\fboxsep}{0pt}
    \fbox{\includegraphics[width=0.25\linewidth]{mobile-report-detail-screen.png}}
    \caption{Экран карточки обращения}
    \label{fig:impl-detail}
\end{figure}

\textbf{Экран уведомлений} (рисунок~\ref{fig:impl-notif}) предназначен для информирования пользователя о смене статусов его обращений. Реализован в виде списка уведомлений с иконками типа события, текстовым содержанием и временной меткой. Реализован индикатор непрочитанных уведомлений.

\begin{figure}[H]
    \centering
    \setlength{\fboxrule}{0.4pt}
    \setlength{\fboxsep}{0pt}
    \fbox{\includegraphics[width=0.25\linewidth]{mobile-notifications-screen.png}}
    \caption{Экран уведомлений}
    \label{fig:impl-notif}
\end{figure}

\textbf{Экран профиля} (рисунок~\ref{fig:impl-profile}) предоставляет доступ к личным данным и настройкам. Включает аватар, имя, блок статистики (общее количество обращений, выполненных и в работе), навигационное меню и кнопку выхода из учётной записи.

\begin{figure}[H]
    \centering
    \setlength{\fboxrule}{0.4pt}
    \setlength{\fboxsep}{0pt}
    \fbox{\includegraphics[width=0.25\linewidth]{mobile-profile-screen.png}}
    \caption{Экран профиля пользователя}
    \label{fig:impl-profile}
\end{figure}

Для обеспечения единообразия интерфейса созданы переиспользуемые компоненты: Button, Input, Card, LoadingIndicator, ErrorDisplay, ReportCard, ImageCarousel. Компоненты применяются на различных экранах и обеспечивают согласованный визуальный стиль. Реализована адаптивность интерфейса с использованием flexbox-разметки, относительных размеров и хука useWindowDimensions.

\subsection{Ключевые алгоритмы мобильного приложения}

\subsubsection{Создание обращения с проверкой дубликатов}

Алгоритм создания обращения реализует сценарий, описанный на диаграмме последовательности (рисунок~\ref{fig:seq-create}). Клиентская часть выполняет следующую последовательность действий:
\begin{enumerate}
    \item пользователь заполняет форму: выбирает категорию, вводит описание, прикрепляет фотографию и указывает местоположение на карте или подтверждает автоматически определённые координаты;
    \item приложение выполняет первичную валидацию: проверка заполненности обязательных полей, наличие фотографии, корректность координат;
    \item при нажатии кнопки <<Отправить>> приложение отправляет запрос POST /api/reports на сервер;
    \item при получении ответа о наличии дубликатов приложение отображает модальное окно с предложением подтвердить существующее обращение или создать новое;
    \item при выборе подтверждения отправляется запрос POST /api/reports/\{id\}/confirm;
    \item при создании нового обращения сервер возвращает созданную сущность со статусом New; приложение обновляет список обращений и карту.
\end{enumerate}

Обработка ошибок сети и серверной валидации реализована через централизованный обработчик в API-клиенте с отображением сообщений пользователю.

\subsubsection{Работа с геолокацией}

Алгоритм работы с геолокацией обеспечивает определение координат для предзаполнения местоположения при создании обращения и отображения текущей позиции на карте:
\begin{enumerate}
    \item при первом обращении к функционалу геолокации приложение запрашивает разрешение через API expo-location;
    \item при предоставлении разрешения выполняется запрос getCurrentPositionAsync с параметрами точности и таймаута;
    \item полученные координаты передаются в компонент карты для отображения маркера текущего местоположения и предзаполнения поля в форме;
    \item при отказе в разрешении предлагается ручное указание на карте;
    \item реализована проверка доступности служб геолокации с информированием пользователя при их отключении.
\end{enumerate}

\subsubsection{Управление аутентификацией}

Алгоритм управления аутентификацией обеспечивает проверку учётных данных, сохранение сессии и доступ к защищённым разделам. На рисунке~\ref{fig:alg-auth} представлена блок-схема алгоритма.

\begin{figure}[H]
    \centering
    \setlength{\fboxrule}{0.4pt}
    \setlength{\fboxsep}{0pt}
    \fbox{\includegraphics[width=0.35\textwidth]{alg-auth.png}}
    \caption{Алгоритм управления аутентификацией}
    \label{fig:alg-auth}
\end{figure}

\begin{enumerate}
    \item пользователь вводит адрес электронной почты и пароль в форме входа;
    \item клиентское приложение отправляет данные на сервер для проверки учётных данных;
    \item при успешной аутентификации сервер возвращает данные сессии, которые сохраняются в защищённом хранилище expo-secure-store;
    \item при обращении к защищённым разделам токен автоматически добавляется к HTTP-запросам через перехватчик Axios;
    \item при получении ответа 401 (Unauthorized) выполняется автоматическое обновление токена и повторный запрос;
    \item при запуске приложения выполняется попытка восстановления ранее сохранённой сессии;
    \item при завершении сеанса локальные данные аутентификации удаляются.
\end{enumerate}
